%------------------------------------------------------------------------------%
\begin{question}
  Calcular a ``área'' sob o gráfico de \(f(x)=x^2\) entre 0 e $b$
\end{question}

%------------------------------------------------------------------------------%
\begin{resolution}{Solução}
  A ``área'' sob o gráfico é
  \[
    \pause
    A
    \;=\; \int_a^b f(x) \,dx
    \pause
    \;=\; \int_0^b x^2  \,dx
  \]
\end{resolution}

%------------------------------------------------------------------------------%
\begin{resolution}{Solução}
  Calculamos a primitiva de $f(x)=x^2$
  \[
    F(x)
    \;=\; \int x^2 \,dx
    \pause
    \;=\; \frac{x^{2+1}}{2+1} + C
    \pause
    \;=\; \frac{x^3}{3} + C
  \]
\end{resolution}

%------------------------------------------------------------------------------%
\begin{resolution}{Solução}
  Pelo Teorema Fundamental do Cálculo,
  \[
    A
    \;=\; F(x) \evalat{0}{b}
    \pause
    \;=\; F(b) - F(0)
    \pause
    \;=\; \left[ \frac{b^3}{3} + C \right]
      -   \left[ \frac{0^3}{3} + C \right]
    \pause
    \;=\; \frac{b^3}{3}
  \]
\end{resolution}

%------------------------------------------------------------------------------%